\chapter{Fix 42: Thermodynamic Stability via the Feshbach-Schur Map}
\label{ch:fix42_thermodynamic_stability}

\section{Context}
Proving a gap $\Delta > 0$ on a finite lattice is insufficient; we must prove $\liminf_{L \to \infty} \Delta_L > 0$. Boundary conditions could theoretically introduce low-energy states that close the gap. We resolve this using \textbf{Isospectral Renormalization} (Appendix R.6 and Appendix AV).

\section{Theorem H.1 (Gap Superadditivity)}
Let $H_\Lambda$ be the Hamiltonian on a large domain $\Lambda$. We can decompose the system into decoupled blocks with Hamiltonian $H_0 = \sum H_{\text{block}}$ and interaction $V$. The resolvent $R(z) = (H_\Lambda - z)^{-1}$ satisfies the Combes-Thomas bound:
\begin{equation}
\| e^{\alpha |x|} R(z) e^{-\alpha |x|} \| \le \frac{C}{\text{dist}(z, \sigma(H))}
\end{equation}
This ensures that the interaction energy between distant vacuum fluctuations decays exponentially, preventing the gap from collapsing.

\section{Proof Derivation}
\subsection{1. Feshbach-Schur Projection}
Let $P$ be the projection onto the low-energy subspace of the decoupled blocks (energy $E < \lambda$) and $\bar{P} = 1 - P$. The effective Hamiltonian on the $P$ subspace is:
\begin{equation}
H_{eff}(z) = PHP - PH\bar{P} (\bar{P}H\bar{P} - z)^{-1} \bar{P}HP
\end{equation}.

\subsection{2. Neumann Series Convergence}
Since $\bar{P}$ projects onto high-energy states (kinetic energy $\gg$ interaction), the resolvent $(\bar{P}H\bar{P} - z)^{-1}$ can be expanded in a convergent Neumann series. The norm is bounded by $1/(\text{Gap}_0 - \|V\|)$.

\subsection{3. Locality of $H_{eff}$}
The Feshbach map preserves the locality of the interactions up to exponential tails. The interaction between two blocks $B_i$ and $B_j$ in $H_{eff}$ decays as $e^{-m \text{dist}(i,j)}$.

\subsection{4. Gap Stability}
Since the effective interactions between vacuum sectors of different blocks are exponentially suppressed, the splitting of the ground state energy (tunneling) is exponentially small. The gap to the first excited state (a particle in one block) remains bounded away from zero uniformly in $L$.
